\documentclass[10pt,a4paper]{amsart}
\usepackage[margin=19mm]{geometry}
\usepackage{amsmath,amssymb,mathtools}
\usepackage[scr=rsfs,cal=euler]{mathalfa}
\usepackage{xfrac}
\usepackage[unicode,hidelinks,bookmarks=false]{hyperref}
\newcommand{\ie}{i.e.\ }
\newcommand{\cf}{cf.\ }
\newcommand{\resp}{respectively\ }
\newcommand{\Subsection}{\subsection}
\providecommand{\nobreakdash}{\nobreak\hbox{-}}
\setlength{\emergencystretch}{2em}
\multlinegap0pt
\usepackage{mathtools}
\usepackage{tensor}
\usepackage{bm}
\usepackage{tikz}
\usetikzlibrary{matrix,arrows,decorations.pathmorphing}
\usepackage{tikz-cd}
\usepackage{eucal}
\usepackage{mathrsfs}
\usepackage{dsfont}
\usepackage{enumitem}
\usepackage{stmaryrd}
\newtheorem{thm}{Theorem}[section]
\newtheorem{prop}[thm]{Proposition}
\newtheorem{lmm}[thm]{Lemma}
\newtheorem{cor}[thm]{Corollary}
\newtheorem{dfn}[thm]{Definition}
\newtheorem{rmk}[thm]{Remark}
\newtheorem{ex}[thm]{Example}
\newtheorem{notat}{Notation}
\DeclareMathOperator{\Aut}{Aut}
\DeclareMathOperator{\coeq}{coeq}
\DeclareMathOperator{\coker}{coker}
\DeclareMathOperator{\car}{char}
\DeclareMathOperator{\Ext}{Ext}
\DeclareMathOperator{\Hilb}{Hilb}
\DeclareMathOperator{\Hom}{Hom}
\DeclareMathOperator{\id}{id}
\DeclareMathOperator{\Id}{Id}
\DeclareMathOperator{\Lan}{Lan}
\DeclareMathOperator{\ob}{ob}
\DeclareMathOperator{\Pic}{Pic}
\DeclareMathOperator{\Quot}{Quot}
\DeclareMathOperator{\rk}{rk}
\DeclareMathOperator{\Spec}{Spec}
\DeclareMathOperator{\Supp}{Supp}
\DeclareMathOperator{\Sym}{Sym}
\newcommand{\E}{\mathbb{E}}
\newcommand{\N}{\mathbb{N}}
\newcommand{\Z}{\mathbb{Z}}
\newcommand{\Q}{\mathbb{Q}}
\newcommand{\R}{\mathbb{R}}
\newcommand{\C}{\mathbb{C}}
\newcommand{\K}{\Bbbk}   %arbitrary field
\newcommand{\aff}{\mathbb{A}}   %affine space
\newcommand{\pr}{\mathbb{P}}   %projective space
\newcommand{\shom}{\mathcal{H}\!\mathit{om}}   %hom sheaf
\newcommand{\cO}{\mathcal{O}} %structure sheaf
\newcommand{\fo}{\mathbf{1}}
\newcommand{\z}{\zeta}
\newcommand{\cD}{{\mathcal D}}
\newcommand{\ve}{{\varepsilon}}
\newcommand{\cR}{{\mathcal R}}
\newcommand{\ka}{{\kappa}}
\newcommand{\mb}[1]{\mathbb{#1}}
\newcommand{\mc}[1]{\mathcal{#1}}
\newcommand{\mf}[1]{\mathfrak{#1}}
\newcommand{\mr}[1]{\mathrm{#1}}
\newcommand{\ms}[1]{\mathsf{#1}}
\newcommand{\mx}[1]{\mathscr{#1}}
\newcommand{\iso}{\xrightarrow{\raisebox{-0.8ex}[0ex][0ex]{$\sim$}}}
\newcommand{\into}{\hookrightarrow}
\newcommand{\onto}{\twoheadrightarrow}
\newcommand{\xto}{\xrightarrow}
\newcommand{\1}{\mathds{1}}   %identity object
\newcommand{\op}{\mathrm{op}}   %opposite category
\newcommand{\Stk}{\mathsf{St}_\K}
\newcommand{\Algst}{\mathsf{AlgSt}_\K}
\newcommand{\bal}{\mathcal{B}^\mathrm{bal}_{0,n}(G)}
\newcommand{\balk}{\mathcal{B}^\mathrm{bal}_{0,k}(G)}
\newcommand{\balN}{\mathcal{B}^\mathrm{bal}_{0,N}(G)}
\newcommand{\balNsm}{\mathcal{B}^\mathrm{sm}_{0,N}(G)}
\newcommand{\balsm}{\mathcal{B}^\mathrm{sm}_{0,n}(G)}
\newcommand{\balc}{\mathcal{B}^\mathrm{bal}_{0,\bm{c}}(G)}
\newcommand{\balcsm}{\mathcal{B}^\mathrm{sm}_{0,\bm{c}}(G)}
\newcommand{\balmv}{\mathcal{B}^\mathrm{bal}_{0,\bm{\mf{m}}(v)}(G)}
\newcommand{\balmw}{\mathcal{B}^\mathrm{bal}_{0,\bm{\mf{m}}(w)}(G)}
\newcommand{\balmvsm}{\mathcal{B}^\mathrm{sm}_{0,\bm{\mf{m}}(v)}(G)}
\newcommand{\balmwsm}{\mathcal{B}^\mathrm{sm}_{0,\bm{\mf{m}}(w)}(G)}
\newcommand{\adm}{\mathcal{A}dm_{0,n}(G)}
\newcommand{\Kstk}{K_0(\Algst)}
\newcommand{\Ksnmod}{K_0^{\bm{S}_n}({\Algst})}
\newcommand{\Ksmod}{K_0^{\ms{S}}({\Algst}/ B)}
\newcommand{\KsmodK}{K_0^{\ms{S}}({\Algst})}
\newcommand{\Ksmodr}{K_0^{\ms{S}, r}({\Algst}/B)}
\newcommand{\BG}{{\mathcal{B}G}}
\newcommand{\cyinst}{\overline{\mathcal{I}}_{\bm{\mu}}(\BG)}
\newcommand{\Kcis}{\overline{\mathcal{I}}_\K}
\newcommand{\und}[1]{\underline{#1}}
\newcommand{\syc}{\mathsf{S}_\mathsf{C}}
\begin{document}
\title[Equivariant Cohomological Crepant Resolution Conjecture for ]{Equivariant Cohomological Crepant Resolution Conjecture for ADE-orbifolds}
\author{Qike Li, Fabio Perroni}
\date{}
\maketitle
\noindent\emph{Source and licence.} Equivariant Cohomological Crepant Resolution Conjecture for ADE-orbifolds. Version arXiv:2609.17494v1 (CC BY 4.0). This material is distributed under the Creative Commons Attribution 4.0 International licence, http://creativecommons.org/licenses/by/4.0/, as recorded in the arXiv OAI licence metadata for this exact version and shown on the abstract page https://arxiv.org/abs/2609.17494v1. Full rights notice: licenses/arxiv-2609.17494-CC-BY-4.0.txt.\par\smallskip
\noindent\textit{Editorial scope.} Counted unit: the original \texttt{thm} environment \texttt{-} at source lines 673--675 together with its complete proof at lines 676--701; the delivered document keeps source lines 486--702 so that every referenced label and every citation resolves inside it. Native editable source is the paper's own concatenated TeX; the AMS wrapper is editorial formatting only. No publisher class, upstream script or unrelated artwork is redistributed.
\subsection{Conventions}\label{conventions}
We work over the field of complex numbers $\C$.
In particular, by a variety we mean a complex algebraic variety and vector bundles will be complex. 
A vector bundle on a disjoint union 
of varieties is given by a vector bundle on each connected component, possibly having different ranks on each 
connected component; its top Chern class is the cohomology class restricting to the top Chern class on each 
connected component. 
All group actions are left group actions.

For a topological space $Y$, 
$H^*(Y)=\oplus_i H^i(Y, \Q)$ denotes the graded singular cohomology group of $Y$ with rational coefficients. 
If we don't take into account the graded structure, we write simply $H(Y)$. When $Y$ is a variety, $H(Y)$
is the cohomology of the complex analytic space associated to $Y$. 

If $Y$ is a differentiable manifold and $T$ is a Lie group acting on it by diffeomorphisms, 
we denote by $H_T(Y)$ the $T$-equivariant cohomology of $Y$,
either defined via the Borel construction as the cohomology of $Y\times_T \E$, 
or as the de Rham cohomology of the stack $[Y/T]$ (\cite{Beh04}, \cite[Ch. 2]{BGNX}). 
In all the applications we deal with in this article, $Y$ is a complex algebraic variety and $T$ is a complex linear algebraic group
acting algebraically on $Y$. Then $H^i_T(Y)$ coincides with the equivariant cohomology group of $Y$, as defined in 
\cite{AF}, for $i\geq 0$. The reader may freely refer to this latter reference regarding equivariant cohomology.
For the sake of clarity, we should mention that, in the Borel construction $\E$ is a contractible topological space
with a free $T$-action. While, in the algebraic framework of \cite{AF} (compare also \cite{EG}) $\E$ is a nonsingular 
connected algebraic variety (hence finite dimensional) such that $H^i(\E)=0$ for every $0<i<N$ (for $N$ a positive integer or infinity),
and with a free $T$-action such that $\E \to \E/T$ is a principal $T$-bundle locally trivial in the complex topology.
Then $H_T^i(Y) \cong H^i(Y\times_T \E)$, for every $i<N$.



\bigskip

\section{Equivariant Chen-Ruan cohomology}\label{ecrc}
Let $Y$ be a smooth variety, let $G$ be a finite group with the identity $e\in G$ and let $T$ be a complex linear algebraic group 
both acting on $Y$ by regular maps, 
such that these actions commute, i.e. they correspond to a $(G\times T)$-action on $Y$.
Under these hypotheses, we obtain  a strict action by $T$ on the orbifold $[Y/G]$, in the sense of \cite{Romagny}.
In a different context (e.g. when $Y$ is a symplectic manifold and $T, G$ are tori),
the $T$-equivariant Chen-Ruan cohomology of $[Y/G]$, that we denote  $H^*_{T, \rm{CR}}([Y/G])$, 
has been defined and studied in \cite{HoMa}. Their approach follows the original definition of \cite{CR04}
and also uses some results from \cite{AGV08}. A similar definition can be found in \cite{Johnson}. 
Since the examples we are interested in don't fit in these frameworks, we review here the
definition of  $H^*_{T, \rm{CR}}([Y/G])$, adapting the presentation of the usual Chen-Ruan cohomology given in \cite{FG03}.

Let 
$$
I_G(Y):= \{ (g, y) \in G\times Y \, | \, gy=y\} \, 
$$
be the inertia variety. Note that $I_G(Y)= \sqcup_{g\in G}Y^g$, where $Y^g= \{ y\in Y \, | \, g y = y\}$
and that it carries a natural action by $T$.

Define the \textit{equivariant inertia cohomology} of $(Y, G)$ by
$$
H_T(Y, G) := H_T(I_G(Y))=  \oplus_{g\in G} H_T(Y^g) \, .
$$
For $\alpha \in H_T(Y^g)$, we denote by $\alpha_g$ the corresponding element 
in the $g$-th summand of $H_T(Y, G)$.

The grading on $H_T(Y, G)$ is defined using the notion of \textit{age}, which we recall now.
In the discussion that follows we consider $Y$ as a complex manifold.
Let $g\in G$ and let $y\in Y^g$. Let $\lambda_1, \ldots , \lambda_{\dim (Y)}$ be the eigenvalues of the 
differential $T_yg$ of the map $g\colon Y \to Y$ at $y$. Since $g$ has finite order,  $\lambda_j = e^{2\pi i r_j}$, 
for unique  $r_j \in \Q \cap [0, 1[$, $j=1, \ldots , \dim (Y)$. 
The age of $g$ in $y$ is defined as $a(g, y):= \sum_{j=1}^{\dim (Y)} r_j$.
By a well known result of Cartan the action of $g$ in a neighborhood of $y$  can be linearized, hence
$a(g, y)$ is constant on the connected component $Z$ that contains $y$. We denote by $a(g, Z)$ the age of $g$ in any point 
$y\in Z$.
Let $\iota \colon Z \hookrightarrow Y$ be the inclusion. Then, for   $\alpha \in H^k_T(Z)$, we assign 
to $(\iota_* \alpha)_g \in H_T(Y, G)$ the degree $k + 2a(g, Z)$. In this way we obtain a rationally graded 
vector space $H_T^*(Y, G)$.


\begin{rmk}
It follows from the definition of age that
$$
a(g, Z) + a(g^{-1}, Z) = \dim (Y) - \dim (Z) = {\rm codim} (Z\subseteq Y) \, .
$$
\end{rmk}
\begin{rmk}
$a(g, Z) \in \Z$ if and only if $\det (T_y g)=1$, for some $y\in Z$.
In particular, $H_T^*(Y, G)$ is integrally graded if the canonical line bundle of $Y$ is $G$-invariant.
\end{rmk}

\begin{notat}
When $Y^g$ is connected (for example if $Y$ is a vector space), we denote $a(g, Y^g)$ by $a(g)$.
\end{notat}

The group $G$ acts on $H_T^*(Y, G)$ as follows: for every $g, h \in G$, $\alpha \in H_T(Y^g)$, define 
$$
h \alpha_g := (h_* \alpha)_{hgh^{-1}} \, .
$$ 
Note that this action preserves the grading.

\begin{dfn}\label{TCR}
Under the previous notation and hypotheses, the $T$-equivariant Chen-Ruan cohomology of $[Y/G]$ is defined as follows:
$$
H^*_{T, {\rm CR}}([Y/G]) := H^*_T(Y, G)^G \, ,
$$
i.e. as the $G$-invariant subspace of $H^*_T(Y, G)$.
\end{dfn}

\begin{rmk}\label{rmkTCR}
Let $\overline{G} \subseteq G$ be a set of representatives of the conjugacy classes of $G$.
Then $H^\ast_{T, {\rm CR}}([Y/G])$ is isomorphic to
$$
\bigoplus_{g\in \overline{G}} H_T^\ast\left(Y^g/C_G(g)\right) \, .
$$ 
where $C_G(g)$ denotes the centralizer of $g$ in $G$.
\end{rmk}

We  now define a bilinear map 
$$
\mu \colon H^*_T(Y, G) \times H^*_T(Y, G) \to H^*_T(Y, G) \, .
$$

Let $g, h \in G$ and let $\langle g, h\rangle \subseteq G$ be the subgroup that they generate.
Let $Y^{g,h} \subseteq Y$ be the complex sub-manifold of points fixed by $\langle g, h\rangle$, that is $Y^{g,h}=Y^g\cap Y^h$.
Note that the inclusions $Y^{g,h} \hookrightarrow Y^g$, $Y^{g,h} \hookrightarrow Y^h$ are $T$-equivariant.
For  $\alpha \in H_T(Y^g)$ and $\beta \in H_T(Y^h)$, we denote by 
$$
\alpha_{|Y^{g,h}} \, , \beta_{|Y^{g,h}} \in H_T(Y^{g,h})
$$
their pull-backs with respect to the above inclusions.

Let $C\to \pr^1$ be the Galois cover, with Galois group $\langle g, h\rangle$, defined in \cite[Appendix]{FG03}
($C$ is denoted $C(\pr^1, g, h, (gh)^{-1}, \langle g, h\rangle)$ in \textit{op. cit.}).
Let us recall that $C\to \pr^1$ is branched over $0, 1, \infty$, and its local monodromy over $0$
(resp. $1, \infty$) is conjugated to $g$ (resp. $h, (gh)^{-1}$). The action of $\langle g, h\rangle$ on $C$ induces a natural action 
on $H^1(C, \cO_C)$. In the  representation ring of $\langle g, h\rangle$ (with complex coefficients) 
we have the following equality \cite[Lemma 8.5]{JKK07}:
\begin{eqnarray}\label{l85JKK07}
H^1(C, \cO_C) = \mathbb C \ominus \C[\langle g, h\rangle] &\oplus& 
\bigoplus_{k=0}^{|g|-1}\frac{k}{|g|} {\rm Ind}^{\langle g, h\rangle}_{\langle g \rangle} \mathbb{V}_{\langle g \rangle, k} \\
&& \bigoplus_{k'=0}^{|h|-1}\frac{k'}{|h|} {\rm Ind}^{\langle g, h\rangle}_{\langle h \rangle} \mathbb{V}_{\langle h \rangle, k'} \nonumber \\
&&  \bigoplus_{k''=0}^{|gh|-1}\frac{k''}{|gh|} {\rm Ind}^{\langle g, h\rangle}_{\langle (gh)^{-1} \rangle} \mathbb{V}_{\langle (gh)^{-1} \rangle, k''} \nonumber \, ,
\end{eqnarray}
where, for every element $x \in G$, $\langle x \rangle$ denotes the subgroup generated by $x$, 
$|x|$ is its order,  $\mathbb{V}_{\langle x \rangle, k}$ is the irreducible representation of $\langle x \rangle$
such that $x$ acts with character $\exp \left( -\frac{2\pi i k}{|x|}\right)$ (for $k=0, \ldots , |x|-1$), and 
${\rm Ind}^{\langle g, h\rangle}_{\langle x \rangle} \mathbb{V}_{\langle x \rangle, k} $ is the induced 
$\langle g, h\rangle$-representation 
$\C[\langle g, h\rangle]\otimes_{\C[\langle x \rangle]} \mathbb{V}_{\langle x \rangle, k}$.

Let us consider the vector bundle $\left( TY_{|Y^{g,h}}\right) \otimes H^1(C, \cO_C)$ on $Y^{g,h}$ with the natural 
$\langle g, h\rangle$-action and define 
$$
F(g,h) := \left[ \left( TY_{|Y^{g,h}}\right) \otimes H^1(C, \cO_C)\right]^{\langle g, h\rangle} \, ,
$$
its $\langle g, h\rangle$-invariant sub-bundle, which is well defined since $\langle g, h\rangle$ is finite. 

Note that $F(g, h)$ is a $T$-equivariant vector bundle on $Y^{g,h}$ (see e.g.   
\cite[Ch. 2.3]{AF} or \cite[Prop. 3.2]{BGNX} for a definition), in fact
$T$ does not act on $H^1(C, \cO_C)$ and $TY_{|Y^{g,h}}$ is a $T$-equivariant vector bundle by standard reasons.
Hence, $F(g, h)$ induces a vector bundle on $Y^{g, h} \times_T \E$, which we denote $F(g, h)_T$. 
By \cite[Lemma 1.12]{FG03} (or \cite[Ch. 8]{JKK07}),
for every connected component $U$ of $Y^{g,h}$, $F(g,h)_{|U}$ and ${F(g,h)_T}_{|U}$ have  rank 
\begin{equation}\label{rkFgh}
{\rm rk}(F(g, h)_{|U}) =a(g, U) + a(h,U) - a(gh, U) - {\rm codim}(U\subseteq Y^{gh}) \, .
\end{equation}
In the following we denote by 
$$
c^T(g, h) \in H_T(Y^{g,h}) 
$$ 
the top equivariant Chern class of $F(g,h)$, which is the top Chern class of $F(g,h)_T$.

\begin{dfn}
The bilinear map 
$$
\mu \colon H^*_T(Y, G) \times H^*_T(Y, G) \to H^*_T(Y, G) \, 
$$
is defined by
$$
\mu (\alpha_g, \beta_h) := \iota_* \left( \alpha_{|Y^{g,h}} \cdot \beta_{|Y^{g,h}} \cdot c^T(g,h) \right) \, ,
$$
for every $g, h \in G$, $\alpha \in H_T(Y^g)$, $\beta \in H_T(Y^h)$; where
$\iota \colon Y^{g, h} \hookrightarrow Y^{gh}$ is the natural  inclusion, $\iota_* \colon H^i_T(Y^{g,h}) \to H^{i+2d}_T(Y^{gh})$ 
is the $T$-equivariant Gysin map ( \cite[Ch. 3.6]{AF}, \cite{BGNX}) and $d= {\rm codim} (Y^{g,h} \subseteq Y^{gh})$.
\end{dfn}

\begin{rmk}\label{mu1}
Note that, if $g=1$, then $C\cong \pr^1$ by the  Riemann-Hurwitz formula, therefore 
$F(1,h)$ has rank $0$ and 
$$
\mu (\alpha_1, \beta_h) =  \left( \alpha_{|Y^{h}} \cdot \beta  \right)_h \, . 
$$
A similar formula holds true if $h=1$.
\end{rmk}


\begin{thm}
The bilinear map $\mu$ defines a graded, $G$-equivariant, associative multiplication on $H^*_T(Y, G)$. 
\end{thm}

\begin{proof}
The theorem follows from the same arguments used in the proof of \cite[Theorem 1.18]{FG03}, suitably adapted
to the $T$-equivariant case.
In particular, the fact that $\mu$ is a graded multiplication is a consequence of the formula \eqref{rkFgh} for the rank of $F(g,h)$.
The $G$-equivariance of $\mu$ follows from the fact that, for $v\in G$, the action map induced by 
$v\colon Y^{g, h} \to Y^{g', h'}$
(where $g'= vgv^{-1}$, $h'= vhv^{-1}$) is $T$-equivariant and $v^* F(g', h') \cong F(g, h)$ by construction. 

To prove the associativity of $\mu$, observe that the excess intersection formula used in \cite{FG03}
holds true for $T$-equivariant cohomology, under our assumptions. In fact, if $S$ is a smooth variety with a $T$-action, 
$S_1$ and $S_2$ are closed smooth $T$-invariant subvarieties, with smooth intersection $U:=S_1 \cap S_2$,   
the excess bundle $E(S, S_1, S_2)$ of $U$ is $T$-equivariant and the bundle induced on $U\times_T \E$
is isomorphic to the excess bundle $E(S\times_T \E, S_1\times_T \E, S_2 \times_T \E)$ of $U\times_T \E$,
where $\E$ is a smooth variety on which $T$ acts freely, satisfying the properties recalled in Section \ref{conventions}
and of sufficiently high dimension. Therefore, for $j_i \colon S_i \hookrightarrow S$
and $\iota_i \colon U \hookrightarrow S_i$ be the natural inclusions, we have that
$$
j_2^*{j_1}_* (\alpha) = {\iota_2}_* \left( e^T(S, S_1, S_2) \cdot \iota_1^*(\alpha) \right) \, 
$$
in the $T$-equivariant cohomology of $S_2$, for every  $\alpha \in H^*_T(S_1)$, where $e^T(S, S_1, S_2)$ is the top equivariant Chern class of $E(S, S_1, S_2)$. 

As a consequence, we have that  Lemma 1.17 of \cite{FG03} holds true in our framework.
Moreover, the vector bundles $F_L$ and $F_R$ in the proof of \cite[Thm. 1.18]{FG03} are isomorphic 
as $T$-equivariant bundles, because $T$ acts trivially on the cohomology of the curve $C$ in \cite[Construction 1.21]{FG03}
(and hence on the cohomology of $\bar{C}$ in \cite[Rem. 1.24]{FG03}). Hence the associativity of $\mu$ follows.
\end{proof}

\begin{thebibliography}{99}
\bibitem[AF24]{AF} David Anderson and William Fulton, \emph{Equivariant cohomology in algebraic geometry}, Cambridge Studies in Advanced Mathematics, vol. 210, Cambridge University Press, Cambridge, 2024. \MR{4655919}
\bibitem[AGV08]{AGV08} D.~Abramovich, T.~Graber, and A.~Vistoli, \emph{Gromov-{W}itten theory of {D}eligne-{M}umford stacks}, Amer. J. Math. \textbf{130} (2008), no.~5, 1337--1398.
\bibitem[BGNX12]{BGNX} Kai Behrend, Gr\'egory Ginot, Behrang Noohi, and Ping Xu, \emph{String topology for stacks}, Ast\'erisque (2012), no.~343, xiv+169. \MR{2977576}
\bibitem[Beh04]{Beh04} K.~Behrend, \emph{Cohomology of stacks. in: Intersection theory and moduli}, ICTP Lecture Notes Series \textbf{19} (2004), 249--294.
\bibitem[CR04]{CR04} Weimin Chen and Yongbin Ruan, \emph{A new cohomology theory of orbifold}, Comm. Math. Phys. \textbf{248} (2004), no.~1, 1--31. \MR{2104605}
\bibitem[EG98]{EG} Dan Edidin and William Graham, \emph{Equivariant intersection theory}, Invent. Math. \textbf{131} (1998), no.~3, 595--634. \MR{1614555}
\bibitem[FG03]{FG03} B.~Fantechi and L.~G\"{o}ttsche, \emph{Orbifold cohomology for global quotients}, Duke Math. J. \textbf{117} (2003), no.~2, 197--227.
\bibitem[HM12]{HoMa} T.~S. Holm and T.~Matsumura, \emph{Equivariant cohomology for {H}amiltonian torus actions on symplectic orbifolds}, Transform. Groups \textbf{17} (2012), no.~3, 717--746. \MR{2956164}
\bibitem[JKK07]{JKK07} Tyler~J. Jarvis, Ralph Kaufmann, and Takashi Kimura, \emph{Stringy {$K$}-theory and the {C}hern character}, Invent. Math. \textbf{168} (2007), no.~1, 23--81. \MR{2285746}
\bibitem[Joh09]{Johnson} Paul~D. Johnson, \emph{Equivariant {G}romov-{W}itten theory of one dimensional stacks}, ProQuest LLC, Ann Arbor, MI, 2009, Thesis (Ph.D.)--University of Michigan. \MR{2713030}
\bibitem[Rom05]{Romagny} Matthieu Romagny, \emph{Group actions on stacks and applications}, Michigan Math. J. \textbf{53} (2005), no.~1, 209--236.
\end{thebibliography}
\end{document}
